\begin{question}[ID=1.1, subtitle={Teorema di Cauchy per $x^3+x$ e $x^2+1$ su $[0,2]$}, tags={cauchy, teoremi}, difficulty=media]
Verifica se le funzioni \(f(x) = x^3 + x\) e \(g(x) = x^2 + 1\) soddisfano le ipotesi del teorema di Cauchy nell'intervallo \([0,\ 2]\) e, in caso affermativo, trova il punto \(c\) che soddisfa il teorema.
\begin{solution}
$\dfrac{f(2)-f(0)}{g(2)-g(0)}=\dfrac{10}{4}=\dfrac{5}{2}$. $\dfrac{3c^2+1}{2c}=\dfrac{5}{2} \Rightarrow 3c^2-5c+1=0 \Rightarrow c=\dfrac{5\pm\sqrt{13}}{6}$
\end{solution}
\end{question}

%-----------------------------%

\begin{question}[ID=1.2, subtitle={Teorema di Cauchy per $\sin(2x)$ e $\cos x$ su $[0,\pi/2]$}, tags={cauchy, teoremi}, difficulty=media]
Verifica se le funzioni \(f(x) = \sin(2x)\) e \(g(x) = \cos(x)\) soddisfano le ipotesi del teorema di Cauchy nell'intervallo \([0;\ \pi/2]\) e, in caso affermativo, trova il punto \(c\) che soddisfa il teorema.
\begin{solution}
$\dfrac{f(\pi/2)-f(0)}{g(\pi/2)-g(0)}=\dfrac{0}{-1}=0$. $\dfrac{2\cos(2c)}{-\sin c}=0 \Rightarrow \cos(2c)=0 \Rightarrow c=\dfrac{\pi}{4}$
\end{solution}
\end{question}

%-----------------------------%

\begin{question}[ID=1.3, subtitle={Teorema di Cauchy per $e^{2x}$ e $e^x$ su $[0,1]$}, tags={cauchy, teoremi}, difficulty=media]
Verifica se le funzioni \(f(x) = e^{2x}\) e \(g(x) = e^x\) soddisfano le ipotesi del teorema di Cauchy nell'intervallo \([0;\ 1]\) e, in caso affermativo, trova il punto \(c\) che soddisfa il teorema.
\begin{solution}
$\dfrac{e^2-1}{e-1}=e+1$. $\dfrac{2e^{2c}}{e^c}=2e^c=e+1 \Rightarrow c=\ln\!\left(\dfrac{e+1}{2}\right)$
\end{solution}
\end{question}

%-----------------------------%

\begin{question}[ID=1.4, subtitle={Teorema di Cauchy per $\ln x$ e $1/x$ su $[1,e]$}, tags={cauchy, teoremi}, difficulty=media]
Date le funzioni \(f(x) = \ln(x)\) e \(g(x) = \dfrac{1}{x}\), verifica se soddisfano le ipotesi del teorema di Cauchy nell'intervallo \([1;\ e]\) e, in caso affermativo, determina i punti che verificano il teorema.
\begin{solution}
$\dfrac{1-0}{1/e-1}=\dfrac{-e}{e-1}$. $\dfrac{1/c}{-1/c^2}=-c=\dfrac{-e}{e-1} \Rightarrow c=\dfrac{e}{e-1}$
\end{solution}
\end{question}

%-----------------------------%

\begin{question}[ID=1.5, subtitle={Teorema di Cauchy per $\arctan x$ e $x$ su $[0,1]$}, tags={cauchy, teoremi}, difficulty=media]
Date le funzioni \(f(x) = \arctan(x)\) e \(g(x) = x\), verifica se soddisfano le ipotesi del teorema di Cauchy nell'intervallo \([0;\ 1]\) e, in caso affermativo, determina i punti che verificano il teorema.
\begin{solution}
$\dfrac{\pi/4}{1}=\dfrac{\pi}{4}$. $\dfrac{1/(1+c^2)}{1}=\dfrac{\pi}{4} \Rightarrow c=\sqrt{\dfrac{4-\pi}{\pi}}$
\end{solution}
\end{question}
